% Options for packages loaded elsewhere
\PassOptionsToPackage{unicode}{hyperref}
\PassOptionsToPackage{hyphens}{url}
\documentclass[
]{article}
\usepackage{xcolor}
\usepackage[margin=1in]{geometry}
\usepackage{amsmath,amssymb}
\setcounter{secnumdepth}{-\maxdimen} % remove section numbering
\usepackage{iftex}
\ifPDFTeX
  \usepackage[T1]{fontenc}
  \usepackage[utf8]{inputenc}
  \usepackage{textcomp} % provide euro and other symbols
\else % if luatex or xetex
  \usepackage{unicode-math} % this also loads fontspec
  \defaultfontfeatures{Scale=MatchLowercase}
  \defaultfontfeatures[\rmfamily]{Ligatures=TeX,Scale=1}
\fi
\usepackage{lmodern}
\ifPDFTeX\else
  % xetex/luatex font selection
\fi
% Use upquote if available, for straight quotes in verbatim environments
\IfFileExists{upquote.sty}{\usepackage{upquote}}{}
\IfFileExists{microtype.sty}{% use microtype if available
  \usepackage[]{microtype}
  \UseMicrotypeSet[protrusion]{basicmath} % disable protrusion for tt fonts
}{}
\makeatletter
\@ifundefined{KOMAClassName}{% if non-KOMA class
  \IfFileExists{parskip.sty}{%
    \usepackage{parskip}
  }{% else
    \setlength{\parindent}{0pt}
    \setlength{\parskip}{6pt plus 2pt minus 1pt}}
}{% if KOMA class
  \KOMAoptions{parskip=half}}
\makeatother
\usepackage{color}
\usepackage{fancyvrb}
\newcommand{\VerbBar}{|}
\newcommand{\VERB}{\Verb[commandchars=\\\{\}]}
\DefineVerbatimEnvironment{Highlighting}{Verbatim}{commandchars=\\\{\}}
% Add ',fontsize=\small' for more characters per line
\usepackage{framed}
\definecolor{shadecolor}{RGB}{248,248,248}
\newenvironment{Shaded}{\begin{snugshade}}{\end{snugshade}}
\newcommand{\AlertTok}[1]{\textcolor[rgb]{0.94,0.16,0.16}{#1}}
\newcommand{\AnnotationTok}[1]{\textcolor[rgb]{0.56,0.35,0.01}{\textbf{\textit{#1}}}}
\newcommand{\AttributeTok}[1]{\textcolor[rgb]{0.13,0.29,0.53}{#1}}
\newcommand{\BaseNTok}[1]{\textcolor[rgb]{0.00,0.00,0.81}{#1}}
\newcommand{\BuiltInTok}[1]{#1}
\newcommand{\CharTok}[1]{\textcolor[rgb]{0.31,0.60,0.02}{#1}}
\newcommand{\CommentTok}[1]{\textcolor[rgb]{0.56,0.35,0.01}{\textit{#1}}}
\newcommand{\CommentVarTok}[1]{\textcolor[rgb]{0.56,0.35,0.01}{\textbf{\textit{#1}}}}
\newcommand{\ConstantTok}[1]{\textcolor[rgb]{0.56,0.35,0.01}{#1}}
\newcommand{\ControlFlowTok}[1]{\textcolor[rgb]{0.13,0.29,0.53}{\textbf{#1}}}
\newcommand{\DataTypeTok}[1]{\textcolor[rgb]{0.13,0.29,0.53}{#1}}
\newcommand{\DecValTok}[1]{\textcolor[rgb]{0.00,0.00,0.81}{#1}}
\newcommand{\DocumentationTok}[1]{\textcolor[rgb]{0.56,0.35,0.01}{\textbf{\textit{#1}}}}
\newcommand{\ErrorTok}[1]{\textcolor[rgb]{0.64,0.00,0.00}{\textbf{#1}}}
\newcommand{\ExtensionTok}[1]{#1}
\newcommand{\FloatTok}[1]{\textcolor[rgb]{0.00,0.00,0.81}{#1}}
\newcommand{\FunctionTok}[1]{\textcolor[rgb]{0.13,0.29,0.53}{\textbf{#1}}}
\newcommand{\ImportTok}[1]{#1}
\newcommand{\InformationTok}[1]{\textcolor[rgb]{0.56,0.35,0.01}{\textbf{\textit{#1}}}}
\newcommand{\KeywordTok}[1]{\textcolor[rgb]{0.13,0.29,0.53}{\textbf{#1}}}
\newcommand{\NormalTok}[1]{#1}
\newcommand{\OperatorTok}[1]{\textcolor[rgb]{0.81,0.36,0.00}{\textbf{#1}}}
\newcommand{\OtherTok}[1]{\textcolor[rgb]{0.56,0.35,0.01}{#1}}
\newcommand{\PreprocessorTok}[1]{\textcolor[rgb]{0.56,0.35,0.01}{\textit{#1}}}
\newcommand{\RegionMarkerTok}[1]{#1}
\newcommand{\SpecialCharTok}[1]{\textcolor[rgb]{0.81,0.36,0.00}{\textbf{#1}}}
\newcommand{\SpecialStringTok}[1]{\textcolor[rgb]{0.31,0.60,0.02}{#1}}
\newcommand{\StringTok}[1]{\textcolor[rgb]{0.31,0.60,0.02}{#1}}
\newcommand{\VariableTok}[1]{\textcolor[rgb]{0.00,0.00,0.00}{#1}}
\newcommand{\VerbatimStringTok}[1]{\textcolor[rgb]{0.31,0.60,0.02}{#1}}
\newcommand{\WarningTok}[1]{\textcolor[rgb]{0.56,0.35,0.01}{\textbf{\textit{#1}}}}
\usepackage{graphicx}
\makeatletter
\newsavebox\pandoc@box
\newcommand*\pandocbounded[1]{% scales image to fit in text height/width
  \sbox\pandoc@box{#1}%
  \Gscale@div\@tempa{\textheight}{\dimexpr\ht\pandoc@box+\dp\pandoc@box\relax}%
  \Gscale@div\@tempb{\linewidth}{\wd\pandoc@box}%
  \ifdim\@tempb\p@<\@tempa\p@\let\@tempa\@tempb\fi% select the smaller of both
  \ifdim\@tempa\p@<\p@\scalebox{\@tempa}{\usebox\pandoc@box}%
  \else\usebox{\pandoc@box}%
  \fi%
}
% Set default figure placement to htbp
\def\fps@figure{htbp}
\makeatother
\setlength{\emergencystretch}{3em} % prevent overfull lines
\providecommand{\tightlist}{%
  \setlength{\itemsep}{0pt}\setlength{\parskip}{0pt}}
\usepackage{bookmark}
\IfFileExists{xurl.sty}{\usepackage{xurl}}{} % add URL line breaks if available
\urlstyle{same}
\hypersetup{
  hidelinks,
  pdfcreator={LaTeX via pandoc}}

\author{}
\date{\vspace{-2.5em}}

\begin{document}

\section{Nash Equilibrium}\label{sec-nash-equilibrium}

\begin{quote}
Definition, existence, computation, and multiplicity of Nash equilibrium
--- the central solution concept in non-cooperative game theory.
\end{quote}

\subsection*{Learning objectives}\label{learning-objectives}
\addcontentsline{toc}{subsection}{Learning objectives}

\begin{itemize}
\tightlist
\item
  Define Nash equilibrium in terms of best responses and unilateral
  deviations.
\item
  State Nash's existence theorem and understand why mixed strategies
  guarantee existence.
\item
  Compute pure-strategy and mixed-strategy Nash equilibria for 2×2 games
  in R.
\item
  Visualize best-response correspondences and identify equilibria
  graphically.
\end{itemize}

\subsection{Motivation}\label{motivation}

Consider two competing firms choosing prices simultaneously. Each firm
wants to maximize profit, but profit depends on the rival's price. Firm
A cannot simply optimize in isolation --- it must reason about Firm B's
choice, knowing that Firm B is reasoning about Firm A in exactly the
same way.

This mutual consistency requirement is the heart of Nash equilibrium: a
strategy profile where no player can improve their payoff by changing
their strategy alone, given what everyone else is doing. It is the most
widely used solution concept in non-cooperative game theory, and the
starting point for almost every applied analysis of strategic
interaction.

\subsection{Theory}\label{theory}

\subsubsection{Best responses}\label{best-responses}

Given a finite game \(G = (N, (A_i)_{i \in N}, (u_i)_{i \in N})\) with
player set \(N\), action sets \(A_i\), and payoff functions \(u_i\),
player \(i\)'s \textbf{best response} to a strategy profile \(a_{-i}\)
of the other players is any action that maximizes \(i\)'s payoff:

\begin{equation}
BR_i(a_{-i}) = \arg\max_{a_i \in A_i} u_i(a_i, a_{-i})
(\#eq:best-response)
\end{equation}

\subsubsection{Nash equilibrium defined}\label{nash-equilibrium-defined}

\textbf{Definition: Nash Equilibrium}

A strategy profile \(a^* = (a_1^*, \ldots, a_n^*)\) is a \textbf{Nash
equilibrium} if every player's strategy is a best response to the
others:

\begin{equation}
u_i(a_i^*, a_{-i}^*) \geq u_i(a_i, a_{-i}^*) \quad \forall\, a_i \in A_i, \; \forall\, i \in N
(\#eq:nash-def)
\end{equation}

Equivalently, \(a_i^* \in BR_i(a_{-i}^*)\) for all \(i\).

The definition captures a stability property: at a Nash equilibrium, no
player has a profitable \textbf{unilateral} deviation. This does not
mean the outcome is socially optimal (the Prisoner's Dilemma shows
otherwise), nor that players coordinate easily --- it means the profile
is self-enforcing once reached.

\subsubsection{Existence}\label{existence}

Nash {[}-@nash1950{]} proved that every finite game --- one with
finitely many players and finitely many actions per player --- has at
least one Nash equilibrium, possibly in mixed strategies. The proof uses
Kakutani's fixed-point theorem applied to the best-response
correspondence.

\textbf{Theorem: Nash's Existence Theorem}

Every finite game has at least one Nash equilibrium in mixed strategies.

The theorem does not guarantee a \emph{pure-strategy} equilibrium. The
game Matching Pennies, for instance, has no pure-strategy equilibrium
but has a unique mixed equilibrium where each player randomizes 50-50
(see @ref(sec-mixed-strategies)).

\subsubsection{Multiplicity}\label{multiplicity}

Games can have zero, one, or many pure-strategy equilibria:

\begin{itemize}
\tightlist
\item
  \textbf{Matching Pennies} --- zero pure, one mixed.
\item
  \textbf{Battle of the Sexes} --- two pure, one mixed.
\item
  \textbf{Prisoner's Dilemma} --- one pure (Defect, Defect), which is
  also the unique NE overall.
\end{itemize}

When multiple equilibria exist, selecting among them is a separate
problem --- the \textbf{equilibrium selection} problem --- addressed by
refinements such as subgame perfection (@ref(sec-extensive-form)),
trembling-hand perfection, and risk dominance.

\subsection{Implementation in R}\label{sec-nash-implementation}

We implement two approaches: a pure-R solver for 2×2 games and a
visualization of best-response correspondences.

\subsubsection{Finding pure-strategy Nash
equilibria}\label{finding-pure-strategy-nash-equilibria}

\begin{Shaded}
\begin{Highlighting}[]
\FunctionTok{source}\NormalTok{(here}\SpecialCharTok{::}\FunctionTok{here}\NormalTok{(}\StringTok{"R"}\NormalTok{, }\StringTok{"solvers.R"}\NormalTok{))}

\CommentTok{\# Prisoner\textquotesingle{}s Dilemma payoff matrices}
\CommentTok{\# Rows = Player 1 actions (C, D), Cols = Player 2 actions (C, D)}
\NormalTok{A }\OtherTok{\textless{}{-}} \FunctionTok{matrix}\NormalTok{(}\FunctionTok{c}\NormalTok{(}\DecValTok{3}\NormalTok{, }\DecValTok{0}\NormalTok{, }\DecValTok{5}\NormalTok{, }\DecValTok{1}\NormalTok{), }\AttributeTok{nrow =} \DecValTok{2}\NormalTok{, }\AttributeTok{byrow =} \ConstantTok{TRUE}\NormalTok{,}
            \AttributeTok{dimnames =} \FunctionTok{list}\NormalTok{(}\FunctionTok{c}\NormalTok{(}\StringTok{"C"}\NormalTok{, }\StringTok{"D"}\NormalTok{), }\FunctionTok{c}\NormalTok{(}\StringTok{"C"}\NormalTok{, }\StringTok{"D"}\NormalTok{)))}
\NormalTok{B }\OtherTok{\textless{}{-}} \FunctionTok{matrix}\NormalTok{(}\FunctionTok{c}\NormalTok{(}\DecValTok{3}\NormalTok{, }\DecValTok{5}\NormalTok{, }\DecValTok{0}\NormalTok{, }\DecValTok{1}\NormalTok{), }\AttributeTok{nrow =} \DecValTok{2}\NormalTok{, }\AttributeTok{byrow =} \ConstantTok{TRUE}\NormalTok{,}
            \AttributeTok{dimnames =} \FunctionTok{list}\NormalTok{(}\FunctionTok{c}\NormalTok{(}\StringTok{"C"}\NormalTok{, }\StringTok{"D"}\NormalTok{), }\FunctionTok{c}\NormalTok{(}\StringTok{"C"}\NormalTok{, }\StringTok{"D"}\NormalTok{)))}

\FunctionTok{cat}\NormalTok{(}\StringTok{"Player 1 payoffs:}\SpecialCharTok{\textbackslash{}n}\StringTok{"}\NormalTok{)}
\end{Highlighting}
\end{Shaded}

\begin{verbatim}
#> Player 1 payoffs:
\end{verbatim}

\begin{Shaded}
\begin{Highlighting}[]
\NormalTok{A}
\end{Highlighting}
\end{Shaded}

\begin{verbatim}
#>   C D
#> C 3 0
#> D 5 1
\end{verbatim}

\begin{Shaded}
\begin{Highlighting}[]
\FunctionTok{cat}\NormalTok{(}\StringTok{"}\SpecialCharTok{\textbackslash{}n}\StringTok{Player 2 payoffs:}\SpecialCharTok{\textbackslash{}n}\StringTok{"}\NormalTok{)}
\end{Highlighting}
\end{Shaded}

\begin{verbatim}
#> 
#> Player 2 payoffs:
\end{verbatim}

\begin{Shaded}
\begin{Highlighting}[]
\NormalTok{B}
\end{Highlighting}
\end{Shaded}

\begin{verbatim}
#>   C D
#> C 3 5
#> D 0 1
\end{verbatim}

\begin{Shaded}
\begin{Highlighting}[]
\NormalTok{pure\_ne }\OtherTok{\textless{}{-}} \FunctionTok{solve\_2x2\_pure\_nash}\NormalTok{(A, B)}
\FunctionTok{cat}\NormalTok{(}\StringTok{"}\SpecialCharTok{\textbackslash{}n}\StringTok{Pure{-}strategy Nash equilibria:}\SpecialCharTok{\textbackslash{}n}\StringTok{"}\NormalTok{)}
\end{Highlighting}
\end{Shaded}

\begin{verbatim}
#> 
#> Pure-strategy Nash equilibria:
\end{verbatim}

\begin{Shaded}
\begin{Highlighting}[]
\ControlFlowTok{for}\NormalTok{ (eq }\ControlFlowTok{in}\NormalTok{ pure\_ne) \{}
  \FunctionTok{cat}\NormalTok{(}\FunctionTok{sprintf}\NormalTok{(}\StringTok{"  (\%s, \%s) with payoffs (\%d, \%d)}\SpecialCharTok{\textbackslash{}n}\StringTok{"}\NormalTok{,}
              \FunctionTok{rownames}\NormalTok{(A)[eq[}\DecValTok{1}\NormalTok{]], }\FunctionTok{colnames}\NormalTok{(A)[eq[}\DecValTok{2}\NormalTok{]],}
\NormalTok{              A[eq[}\DecValTok{1}\NormalTok{], eq[}\DecValTok{2}\NormalTok{]], B[eq[}\DecValTok{1}\NormalTok{], eq[}\DecValTok{2}\NormalTok{]]))}
\NormalTok{\}}
\end{Highlighting}
\end{Shaded}

\begin{verbatim}
#>   (D, D) with payoffs (1, 1)
\end{verbatim}

\subsubsection{Battle of the Sexes --- multiple
equilibria}\label{battle-of-the-sexes-multiple-equilibria}

\begin{Shaded}
\begin{Highlighting}[]
\CommentTok{\# Battle of the Sexes}
\NormalTok{A\_bos }\OtherTok{\textless{}{-}} \FunctionTok{matrix}\NormalTok{(}\FunctionTok{c}\NormalTok{(}\DecValTok{3}\NormalTok{, }\DecValTok{0}\NormalTok{, }\DecValTok{0}\NormalTok{, }\DecValTok{2}\NormalTok{), }\AttributeTok{nrow =} \DecValTok{2}\NormalTok{, }\AttributeTok{byrow =} \ConstantTok{TRUE}\NormalTok{,}
                \AttributeTok{dimnames =} \FunctionTok{list}\NormalTok{(}\FunctionTok{c}\NormalTok{(}\StringTok{"Opera"}\NormalTok{, }\StringTok{"Football"}\NormalTok{), }\FunctionTok{c}\NormalTok{(}\StringTok{"Opera"}\NormalTok{, }\StringTok{"Football"}\NormalTok{)))}
\NormalTok{B\_bos }\OtherTok{\textless{}{-}} \FunctionTok{matrix}\NormalTok{(}\FunctionTok{c}\NormalTok{(}\DecValTok{2}\NormalTok{, }\DecValTok{0}\NormalTok{, }\DecValTok{0}\NormalTok{, }\DecValTok{3}\NormalTok{), }\AttributeTok{nrow =} \DecValTok{2}\NormalTok{, }\AttributeTok{byrow =} \ConstantTok{TRUE}\NormalTok{,}
                \AttributeTok{dimnames =} \FunctionTok{list}\NormalTok{(}\FunctionTok{c}\NormalTok{(}\StringTok{"Opera"}\NormalTok{, }\StringTok{"Football"}\NormalTok{), }\FunctionTok{c}\NormalTok{(}\StringTok{"Opera"}\NormalTok{, }\StringTok{"Football"}\NormalTok{)))}

\FunctionTok{cat}\NormalTok{(}\StringTok{"Battle of the Sexes:}\SpecialCharTok{\textbackslash{}n}\StringTok{"}\NormalTok{)}
\end{Highlighting}
\end{Shaded}

\begin{verbatim}
#> Battle of the Sexes:
\end{verbatim}

\begin{Shaded}
\begin{Highlighting}[]
\FunctionTok{cat}\NormalTok{(}\StringTok{"Player 1 payoffs:}\SpecialCharTok{\textbackslash{}n}\StringTok{"}\NormalTok{)}
\end{Highlighting}
\end{Shaded}

\begin{verbatim}
#> Player 1 payoffs:
\end{verbatim}

\begin{Shaded}
\begin{Highlighting}[]
\NormalTok{A\_bos}
\end{Highlighting}
\end{Shaded}

\begin{verbatim}
#>          Opera Football
#> Opera        3        0
#> Football     0        2
\end{verbatim}

\begin{Shaded}
\begin{Highlighting}[]
\FunctionTok{cat}\NormalTok{(}\StringTok{"}\SpecialCharTok{\textbackslash{}n}\StringTok{Player 2 payoffs:}\SpecialCharTok{\textbackslash{}n}\StringTok{"}\NormalTok{)}
\end{Highlighting}
\end{Shaded}

\begin{verbatim}
#> 
#> Player 2 payoffs:
\end{verbatim}

\begin{Shaded}
\begin{Highlighting}[]
\NormalTok{B\_bos}
\end{Highlighting}
\end{Shaded}

\begin{verbatim}
#>          Opera Football
#> Opera        2        0
#> Football     0        3
\end{verbatim}

\begin{Shaded}
\begin{Highlighting}[]
\NormalTok{pure\_bos }\OtherTok{\textless{}{-}} \FunctionTok{solve\_2x2\_pure\_nash}\NormalTok{(A\_bos, B\_bos)}
\NormalTok{mixed\_bos }\OtherTok{\textless{}{-}} \FunctionTok{solve\_2x2\_mixed\_nash}\NormalTok{(A\_bos, B\_bos)}

\FunctionTok{cat}\NormalTok{(}\StringTok{"}\SpecialCharTok{\textbackslash{}n}\StringTok{Pure{-}strategy NE:}\SpecialCharTok{\textbackslash{}n}\StringTok{"}\NormalTok{)}
\end{Highlighting}
\end{Shaded}

\begin{verbatim}
#> 
#> Pure-strategy NE:
\end{verbatim}

\begin{Shaded}
\begin{Highlighting}[]
\ControlFlowTok{for}\NormalTok{ (eq }\ControlFlowTok{in}\NormalTok{ pure\_bos) \{}
  \FunctionTok{cat}\NormalTok{(}\FunctionTok{sprintf}\NormalTok{(}\StringTok{"  (\%s, \%s)}\SpecialCharTok{\textbackslash{}n}\StringTok{"}\NormalTok{, }\FunctionTok{rownames}\NormalTok{(A\_bos)[eq[}\DecValTok{1}\NormalTok{]], }\FunctionTok{colnames}\NormalTok{(A\_bos)[eq[}\DecValTok{2}\NormalTok{]]))}
\NormalTok{\}}
\end{Highlighting}
\end{Shaded}

\begin{verbatim}
#>   (Opera, Opera)
#>   (Football, Football)
\end{verbatim}

\begin{Shaded}
\begin{Highlighting}[]
\FunctionTok{cat}\NormalTok{(}\FunctionTok{sprintf}\NormalTok{(}\StringTok{"}\SpecialCharTok{\textbackslash{}n}\StringTok{Mixed{-}strategy NE: p = \%.2f, q = \%.2f}\SpecialCharTok{\textbackslash{}n}\StringTok{"}\NormalTok{,}
\NormalTok{            mixed\_bos}\SpecialCharTok{$}\NormalTok{p, mixed\_bos}\SpecialCharTok{$}\NormalTok{q))}
\end{Highlighting}
\end{Shaded}

\begin{verbatim}
#> 
#> Mixed-strategy NE: p = 0.60, q = 0.40
\end{verbatim}

\subsubsection{Best-response correspondence
plot}\label{best-response-correspondence-plot}

\begin{Shaded}
\begin{Highlighting}[]
\CommentTok{\# For a 2x2 game, let p = P(Player 1 plays row 1), q = P(Player 2 plays col 1)}
\CommentTok{\# Player 2\textquotesingle{}s expected payoff from col 1: p * B[1,1] + (1{-}p) * B[2,1]}
\CommentTok{\# Player 2\textquotesingle{}s expected payoff from col 2: p * B[1,2] + (1{-}p) * B[2,2]}
\CommentTok{\# Player 2 is indifferent when these are equal}

\NormalTok{p\_seq }\OtherTok{\textless{}{-}} \FunctionTok{seq}\NormalTok{(}\DecValTok{0}\NormalTok{, }\DecValTok{1}\NormalTok{, }\AttributeTok{length.out =} \DecValTok{200}\NormalTok{)}

\CommentTok{\# Player 1\textquotesingle{}s best response: BR1(q)}
\CommentTok{\# EU1(row1) = q*A[1,1] + (1{-}q)*A[1,2] = q*3}
\CommentTok{\# EU1(row2) = q*A[2,1] + (1{-}q)*A[2,2] = (1{-}q)*2}
\CommentTok{\# Indifferent when 3q = 2(1{-}q) =\textgreater{} q = 2/5}
\NormalTok{br1\_q }\OtherTok{\textless{}{-}} \DecValTok{2} \SpecialCharTok{/}\NormalTok{ (}\DecValTok{3} \SpecialCharTok{+} \DecValTok{2}\NormalTok{)  }\CommentTok{\# q at which P1 is indifferent}

\CommentTok{\# Player 2\textquotesingle{}s best response: BR2(p)}
\CommentTok{\# EU2(col1) = p*B[1,1] + (1{-}p)*B[2,1] = 2p}
\CommentTok{\# EU2(col2) = p*B[1,2] + (1{-}p)*B[2,2] = 3(1{-}p)}
\CommentTok{\# Indifferent when 2p = 3(1{-}p) =\textgreater{} p = 3/5}
\NormalTok{br2\_p }\OtherTok{\textless{}{-}} \DecValTok{3} \SpecialCharTok{/}\NormalTok{ (}\DecValTok{2} \SpecialCharTok{+} \DecValTok{3}\NormalTok{)  }\CommentTok{\# p at which P2 is indifferent}

\CommentTok{\# Build best{-}response data}
\NormalTok{br\_data }\OtherTok{\textless{}{-}} \FunctionTok{tibble}\NormalTok{(}
  \AttributeTok{q =} \FunctionTok{c}\NormalTok{(}\DecValTok{0}\NormalTok{, br1\_q, br1\_q, br1\_q, }\DecValTok{1}\NormalTok{),}
  \AttributeTok{p\_br1 =} \FunctionTok{c}\NormalTok{(}\DecValTok{0}\NormalTok{, }\DecValTok{0}\NormalTok{, }\FloatTok{0.5}\NormalTok{, }\DecValTok{1}\NormalTok{, }\DecValTok{1}\NormalTok{),}
  \AttributeTok{type =} \StringTok{"Player 1 BR"}
\NormalTok{)}

\NormalTok{br2\_data }\OtherTok{\textless{}{-}} \FunctionTok{tibble}\NormalTok{(}
  \AttributeTok{p =} \FunctionTok{c}\NormalTok{(}\DecValTok{0}\NormalTok{, br2\_p, br2\_p, br2\_p, }\DecValTok{1}\NormalTok{),}
  \AttributeTok{q\_br2 =} \FunctionTok{c}\NormalTok{(}\DecValTok{0}\NormalTok{, }\DecValTok{0}\NormalTok{, }\FloatTok{0.5}\NormalTok{, }\DecValTok{1}\NormalTok{, }\DecValTok{1}\NormalTok{),}
  \AttributeTok{type =} \StringTok{"Player 2 BR"}
\NormalTok{)}

\CommentTok{\# Equilibrium points}
\NormalTok{eq\_points }\OtherTok{\textless{}{-}} \FunctionTok{tibble}\NormalTok{(}
  \AttributeTok{p =} \FunctionTok{c}\NormalTok{(}\DecValTok{0}\NormalTok{, br2\_p, }\DecValTok{1}\NormalTok{),}
  \AttributeTok{q =} \FunctionTok{c}\NormalTok{(}\DecValTok{0}\NormalTok{, br1\_q, }\DecValTok{1}\NormalTok{),}
  \AttributeTok{label =} \FunctionTok{c}\NormalTok{(}\StringTok{"(Football, Football)"}\NormalTok{, }\StringTok{"Mixed NE"}\NormalTok{, }\StringTok{"(Opera, Opera)"}\NormalTok{)}
\NormalTok{)}

\NormalTok{p\_fig }\OtherTok{\textless{}{-}} \FunctionTok{ggplot}\NormalTok{() }\SpecialCharTok{+}
  \FunctionTok{geom\_path}\NormalTok{(}\AttributeTok{data =}\NormalTok{ br\_data, }\FunctionTok{aes}\NormalTok{(}\AttributeTok{x =}\NormalTok{ q, }\AttributeTok{y =}\NormalTok{ p\_br1),}
            \AttributeTok{colour =}\NormalTok{ okabe\_ito[}\DecValTok{1}\NormalTok{], }\AttributeTok{linewidth =} \FloatTok{1.2}\NormalTok{) }\SpecialCharTok{+}
  \FunctionTok{geom\_path}\NormalTok{(}\AttributeTok{data =}\NormalTok{ br2\_data, }\FunctionTok{aes}\NormalTok{(}\AttributeTok{x =}\NormalTok{ q\_br2, }\AttributeTok{y =}\NormalTok{ p),}
            \AttributeTok{colour =}\NormalTok{ okabe\_ito[}\DecValTok{2}\NormalTok{], }\AttributeTok{linewidth =} \FloatTok{1.2}\NormalTok{) }\SpecialCharTok{+}
  \FunctionTok{geom\_point}\NormalTok{(}\AttributeTok{data =}\NormalTok{ eq\_points, }\FunctionTok{aes}\NormalTok{(}\AttributeTok{x =}\NormalTok{ q, }\AttributeTok{y =}\NormalTok{ p),}
             \AttributeTok{size =} \DecValTok{3}\NormalTok{, }\AttributeTok{colour =}\NormalTok{ okabe\_ito[}\DecValTok{6}\NormalTok{]) }\SpecialCharTok{+}
  \FunctionTok{geom\_text}\NormalTok{(}\AttributeTok{data =}\NormalTok{ eq\_points, }\FunctionTok{aes}\NormalTok{(}\AttributeTok{x =}\NormalTok{ q, }\AttributeTok{y =}\NormalTok{ p, }\AttributeTok{label =}\NormalTok{ label),}
            \AttributeTok{vjust =} \SpecialCharTok{{-}}\DecValTok{1}\NormalTok{, }\AttributeTok{size =} \DecValTok{3}\NormalTok{) }\SpecialCharTok{+}
  \FunctionTok{annotate}\NormalTok{(}\StringTok{"text"}\NormalTok{, }\AttributeTok{x =} \FloatTok{0.85}\NormalTok{, }\AttributeTok{y =} \FloatTok{0.15}\NormalTok{, }\AttributeTok{label =} \StringTok{"BR₁(q)"}\NormalTok{,}
           \AttributeTok{colour =}\NormalTok{ okabe\_ito[}\DecValTok{1}\NormalTok{], }\AttributeTok{fontface =} \StringTok{"bold"}\NormalTok{, }\AttributeTok{size =} \DecValTok{4}\NormalTok{) }\SpecialCharTok{+}
  \FunctionTok{annotate}\NormalTok{(}\StringTok{"text"}\NormalTok{, }\AttributeTok{x =} \FloatTok{0.15}\NormalTok{, }\AttributeTok{y =} \FloatTok{0.85}\NormalTok{, }\AttributeTok{label =} \StringTok{"BR₂(p)"}\NormalTok{,}
           \AttributeTok{colour =}\NormalTok{ okabe\_ito[}\DecValTok{2}\NormalTok{], }\AttributeTok{fontface =} \StringTok{"bold"}\NormalTok{, }\AttributeTok{size =} \DecValTok{4}\NormalTok{) }\SpecialCharTok{+}
  \FunctionTok{scale\_x\_continuous}\NormalTok{(}\AttributeTok{name =} \StringTok{"q  (Player 2: prob. of Opera)"}\NormalTok{,}
                     \AttributeTok{limits =} \FunctionTok{c}\NormalTok{(}\SpecialCharTok{{-}}\FloatTok{0.05}\NormalTok{, }\FloatTok{1.15}\NormalTok{), }\AttributeTok{breaks =} \FunctionTok{seq}\NormalTok{(}\DecValTok{0}\NormalTok{, }\DecValTok{1}\NormalTok{, }\FloatTok{0.2}\NormalTok{)) }\SpecialCharTok{+}
  \FunctionTok{scale\_y\_continuous}\NormalTok{(}\AttributeTok{name =} \StringTok{"p  (Player 1: prob. of Opera)"}\NormalTok{,}
                     \AttributeTok{limits =} \FunctionTok{c}\NormalTok{(}\SpecialCharTok{{-}}\FloatTok{0.05}\NormalTok{, }\FloatTok{1.25}\NormalTok{), }\AttributeTok{breaks =} \FunctionTok{seq}\NormalTok{(}\DecValTok{0}\NormalTok{, }\DecValTok{1}\NormalTok{, }\FloatTok{0.2}\NormalTok{)) }\SpecialCharTok{+}
  \FunctionTok{theme\_publication}\NormalTok{() }\SpecialCharTok{+}
  \FunctionTok{labs}\NormalTok{(}\AttributeTok{title =} \StringTok{"Best{-}Response Correspondences: Battle of the Sexes"}\NormalTok{)}

\NormalTok{p\_fig}
\end{Highlighting}
\end{Shaded}

\begin{figure}

{\centering \includegraphics[width=0.8\linewidth]{/home/runner/work/strategy-in-r/strategy-in-r/docs/pdf-chapters/03-nash-equilibrium_files/figure-latex/best-response-1} 

}

\caption{Best-response correspondences for the Battle of the Sexes. Player 1's best response (orange) and Player 2's best response (blue) intersect at three Nash equilibria: two pure-strategy and one mixed.}\label{fig:best-response}
\end{figure}

\begin{Shaded}
\begin{Highlighting}[]
\FunctionTok{save\_pub\_fig}\NormalTok{(p\_fig, }\StringTok{"best{-}response{-}bos"}\NormalTok{)}
\end{Highlighting}
\end{Shaded}

\begin{verbatim}
#> Warning in grid.Call.graphics(C_text, as.graphicsAnnot(x$label), x$x, x$y, :
#> conversion failure on 'BR₁(q)' in 'mbcsToSbcs': for ₁ (U+2081)
\end{verbatim}

\begin{verbatim}
#> Warning in grid.Call.graphics(C_text, as.graphicsAnnot(x$label), x$x, x$y, :
#> conversion failure on 'BR₂(p)' in 'mbcsToSbcs': for ₂ (U+2082)
\end{verbatim}

\subsection{Worked example}\label{worked-example}

Consider a coordination game where two drivers approach each other on a
road and must choose to swerve Left or Right:

\begin{Shaded}
\begin{Highlighting}[]
\NormalTok{A\_coord }\OtherTok{\textless{}{-}} \FunctionTok{matrix}\NormalTok{(}\FunctionTok{c}\NormalTok{(}\DecValTok{1}\NormalTok{, }\DecValTok{0}\NormalTok{, }\DecValTok{0}\NormalTok{, }\DecValTok{1}\NormalTok{), }\AttributeTok{nrow =} \DecValTok{2}\NormalTok{, }\AttributeTok{byrow =} \ConstantTok{TRUE}\NormalTok{,}
                  \AttributeTok{dimnames =} \FunctionTok{list}\NormalTok{(}\FunctionTok{c}\NormalTok{(}\StringTok{"Left"}\NormalTok{, }\StringTok{"Right"}\NormalTok{), }\FunctionTok{c}\NormalTok{(}\StringTok{"Left"}\NormalTok{, }\StringTok{"Right"}\NormalTok{)))}
\NormalTok{B\_coord }\OtherTok{\textless{}{-}}\NormalTok{ A\_coord  }\CommentTok{\# symmetric game}

\FunctionTok{cat}\NormalTok{(}\StringTok{"Coordination Game (symmetric):}\SpecialCharTok{\textbackslash{}n}\StringTok{"}\NormalTok{)}
\end{Highlighting}
\end{Shaded}

\begin{verbatim}
#> Coordination Game (symmetric):
\end{verbatim}

\begin{Shaded}
\begin{Highlighting}[]
\NormalTok{A\_coord}
\end{Highlighting}
\end{Shaded}

\begin{verbatim}
#>       Left Right
#> Left     1     0
#> Right    0     1
\end{verbatim}

\begin{Shaded}
\begin{Highlighting}[]
\NormalTok{eq\_coord }\OtherTok{\textless{}{-}} \FunctionTok{support\_enumeration}\NormalTok{(A\_coord, B\_coord)}

\FunctionTok{cat}\NormalTok{(}\StringTok{"}\SpecialCharTok{\textbackslash{}n}\StringTok{Pure{-}strategy NE:}\SpecialCharTok{\textbackslash{}n}\StringTok{"}\NormalTok{)}
\end{Highlighting}
\end{Shaded}

\begin{verbatim}
#> 
#> Pure-strategy NE:
\end{verbatim}

\begin{Shaded}
\begin{Highlighting}[]
\ControlFlowTok{for}\NormalTok{ (eq }\ControlFlowTok{in}\NormalTok{ eq\_coord}\SpecialCharTok{$}\NormalTok{pure) \{}
  \FunctionTok{cat}\NormalTok{(}\FunctionTok{sprintf}\NormalTok{(}\StringTok{"  (\%s, \%s)}\SpecialCharTok{\textbackslash{}n}\StringTok{"}\NormalTok{, }\FunctionTok{rownames}\NormalTok{(A\_coord)[eq[}\DecValTok{1}\NormalTok{]], }\FunctionTok{colnames}\NormalTok{(A\_coord)[eq[}\DecValTok{2}\NormalTok{]]))}
\NormalTok{\}}
\end{Highlighting}
\end{Shaded}

\begin{verbatim}
#>   (Left, Left)
#>   (Right, Right)
\end{verbatim}

\begin{Shaded}
\begin{Highlighting}[]
\ControlFlowTok{if}\NormalTok{ (}\SpecialCharTok{!}\FunctionTok{is.null}\NormalTok{(eq\_coord}\SpecialCharTok{$}\NormalTok{mixed)) \{}
  \FunctionTok{cat}\NormalTok{(}\FunctionTok{sprintf}\NormalTok{(}\StringTok{"}\SpecialCharTok{\textbackslash{}n}\StringTok{Mixed NE: p = \%.2f, q = \%.2f}\SpecialCharTok{\textbackslash{}n}\StringTok{"}\NormalTok{,}
\NormalTok{              eq\_coord}\SpecialCharTok{$}\NormalTok{mixed}\SpecialCharTok{$}\NormalTok{p, eq\_coord}\SpecialCharTok{$}\NormalTok{mixed}\SpecialCharTok{$}\NormalTok{q))}
  \FunctionTok{cat}\NormalTok{(}\StringTok{"Expected payoff at mixed NE: "}\NormalTok{,}
\NormalTok{      eq\_coord}\SpecialCharTok{$}\NormalTok{mixed}\SpecialCharTok{$}\NormalTok{p }\SpecialCharTok{*}\NormalTok{ eq\_coord}\SpecialCharTok{$}\NormalTok{mixed}\SpecialCharTok{$}\NormalTok{q }\SpecialCharTok{*} \DecValTok{1} \SpecialCharTok{+}
\NormalTok{      (}\DecValTok{1} \SpecialCharTok{{-}}\NormalTok{ eq\_coord}\SpecialCharTok{$}\NormalTok{mixed}\SpecialCharTok{$}\NormalTok{p) }\SpecialCharTok{*}\NormalTok{ (}\DecValTok{1} \SpecialCharTok{{-}}\NormalTok{ eq\_coord}\SpecialCharTok{$}\NormalTok{mixed}\SpecialCharTok{$}\NormalTok{q) }\SpecialCharTok{*} \DecValTok{1}\NormalTok{, }\StringTok{"}\SpecialCharTok{\textbackslash{}n}\StringTok{"}\NormalTok{)}
\NormalTok{\}}
\end{Highlighting}
\end{Shaded}

\begin{verbatim}
#> 
#> Mixed NE: p = 0.50, q = 0.50
#> Expected payoff at mixed NE:  0.5
\end{verbatim}

The coordination game has two pure NE --- (Left, Left) and (Right,
Right) --- and a mixed NE at \(p = q = 0.5\) with expected payoff 0.5,
which is worse than either pure equilibrium. This illustrates that mixed
equilibria can be Pareto-dominated by pure ones: the randomization
creates a risk of miscoordination.

\subsection{Extensions}\label{extensions}

Nash equilibrium is the foundation on which most of the rest of this
book builds:

\begin{itemize}
\tightlist
\item
  \textbf{Mixed strategies} (@ref(sec-mixed-strategies)) formalize the
  randomization used in the mixed NE above.
\item
  \textbf{Extensive-form games} (@ref(sec-extensive-form)) introduce
  sequential moves and \textbf{subgame perfect equilibrium}, a
  refinement of Nash equilibrium.
\item
  \textbf{Bayesian games} (@ref(sec-bayesian-games)) extend Nash
  equilibrium to settings with incomplete information.
\item
  \textbf{Repeated games} (@ref(sec-repeated-games)) show how the folk
  theorem expands the set of equilibria when games are played
  repeatedly.
\end{itemize}

For a comprehensive treatment of Nash equilibrium and its refinements,
see @osborne2004 (Chapters 2--4).

\subsection*{Exercises}\label{exercises}
\addcontentsline{toc}{subsection}{Exercises}

\begin{enumerate}
\def\labelenumi{\arabic{enumi}.}
\item
  \textbf{Stag Hunt.} Consider the Stag Hunt game with payoffs: both
  hunt Stag yields (4, 4), both hunt Hare yields (3, 3), one hunts Stag
  while the other hunts Hare yields (0, 3) for the Stag hunter and (3,
  0) for the Hare hunter. Find all pure-strategy and mixed-strategy Nash
  equilibria. Which equilibrium is risk-dominant? Which is
  payoff-dominant?
\item
  \textbf{Three equilibria.} Construct a 2×2 game that has exactly two
  pure-strategy Nash equilibria and one mixed-strategy Nash equilibrium.
  Verify your answer using the \texttt{support\_enumeration()} function
  from \texttt{R/solvers.R}.
\item
  \textbf{Best-response plot.} Modify the best-response correspondence
  code in @ref(sec-nash-implementation) to plot the correspondences for
  the game of Chicken (payoffs: (0, 0) for mutual swerve, (−1, 1) and
  (1, −1) for one swerving, (−5, −5) for neither swerving). How many
  equilibria does the plot reveal?
\end{enumerate}

Solutions appear in @ref(sec-solutions).

\end{document}
